%==============================================================
% Advanced Scientific Writing -- Assignment 4
% Research-paper FIRST DRAFT, typeset with the official PSB / World
% Scientific author kit (ws-procs11x85 document class).
%
% Build:  pdflatex twice  (or: latexmk -pdf -cd)
%
% DRAFT: conventional section structure (Introduction, Background,
% Analytical Objective, Experimental Data, Algorithms Used, Metrics,
% Conclusion). Kept deliberately concise (~3 pages). [TODO] markers
% flag passages still to be developed or cited.
%==============================================================

\documentclass{ws-procs11x85}
\usepackage{ws-procs-thm}   % comment out if amsthm/theorem/ntheorem is used

\begin{document}

\title{Nicotinic Acetylcholine Receptor Prediction using VEP for Humans}

\author{Hari Krishna Mohan Raj}

\address{Department of Computational Chemistry,\\
Hochschule Bonn Rhein Sieg, Sankt Augustin, Germany\\
E-mail: harikrishnam0711@gmail.com}

\begin{abstract}
Missense variants in human nicotinic acetylcholine receptor (nAChR)
genes are linked tto congenital myasthenic syndromes, epilepsy, and
related disorders, yet most newly observed variants are of uncertain
significance. Standard variant effect predictors report only whether a
variant is damaging, not whether it causes loss of function (LOF) or
gain of function (GOF) a distinction that is decisive for therapy. We
assemble a curated dataset of experimentally characterised nAChR
missense variants labelled by direction of effect, engineer sequence-
and structure-derived features, and benchmark a panel of machine-learning
classifiers to predict LOF versus GOF. Gradient boosted models perform
best, indicating that direction-aware, receptor-specific prediction is
feasible and can help triage nAChR variants of uncertain significance
toward functional follow up and precision therapy.
\end{abstract}

\keywords{Variant effect prediction; nicotinic acetylcholine receptor;
gain of function; loss of function; ion channel; machine learning.}

% required, do-not-remove
\copyrightinfo{\copyright\ 2026 Hari Krishna Mohan Raj, Karl Kirschner. Open Access chapter
published by World Scientific Publishing Company and distributed under
the terms of the Creative Commons Attribution Non-Commercial (CC
BY-NC) 4.0 License.}

%==============================================================
\section{Introduction}\label{sec:intro}
%==============================================================
Missense variants in the genes encoding nicotinic acetylcholine
receptors (nAChRs) underlie congenital myasthenic syndromes at the
neuromuscular junction and several inherited epilepsies in the brain.
As clinical sequencing becomes routine, new variants are observed far
faster than their functional consequences can be measured, leaving most
as variants of uncertain significance (VUS). Existing computational
predictors such as SIFT, PolyPhen-2,,
CADD, and AlphaMissense report whether a
variant is damaging, but not whether it causes loss of function (LOF) or
gain of function (GOF). For nAChRs this is a critical omission, because
LOF and GOF often demand opposite clinical strategies. This work asks
whether the direction of effect can be predicted computationally, and
treats it as a receptor-specific binary classification problem.

%==============================================================
\section{Background}\label{sec:background}
%==============================================================
nAChRs are pentameric ligand-gated ion channels assembled from a
repertoire of seventeen human subunit genes that combine into muscle-type
and neuronal receptors. The affected subunit largely determines where a
receptor acts and which disorder a variant is associated with, so
subunit identity is as informative as the substitution itself. The
direction of a variant's effect is measurable by electrophysiology, which
reveals whether a mutant receptor conducts too little or too much current,
but such assays are slow and costly and cannot scale to the volume of
variants now being reported. This gap between a clinically decisive
property and the cost of measuring it directly motivates a computational
approach.

[TODO: cite primary genotype-to-phenotype references for the
specific subunit-to-disorder mappings.]

%==============================================================
\section{Analytical Objective}\label{sec:objective}
%==============================================================
The objective is to classify a given nAChR missense variant as LOF or
GOF, rather than merely pathogenic or benign. Unlike general pathogenicity
tools, and unlike direction-aware predictors developed for other channel
families such as funNCion, LoGoFunc, and
PreMode, the model is built specifically around nAChR
biology and subunit identity. The task is framed as supervised binary
classification, evaluated by its ability to recover the experimentally
established direction of effect for held-out variants.

%==============================================================
\section{Experimental Data}\label{sec:data}
%==============================================================
The training data comprise a curated set of experimentally characterised
nAChR missense variants spanning the majority of human subunit genes,
each labelled LOF or GOF from primary electrophysiology
reports. Every variant is represented by an
engineered feature vector combining amino-acid physicochemical properties,
evolutionary substitution scores (BLOSUM62 and Grantham distance),
normalised sequence position, a one-hot encoding of subunit identity, and
structural descriptors derived from experimental cryo-EM structures
(per-residue B-factor, secondary structure, and local packing density).
A majority of variants map successfully onto solved structures, providing
the structural context that sequence-only methods lack. 

[TODO: confirm final dataset size and the exact feature count once the curation pass is
locked.]

%==============================================================
\section{Features, Algorithm and Parameters}\label{sec:algorithms}
%==============================================================
A panel of seven machine learning classifiers was benchmarked on the same
data and features, ranging from linear models and support vector machines
to random forests, gradient-boosted trees (XGBoost and LightGBM), and a
neural network. Models were trained and tuned with nested cross-validation
and Optuna-based hyperparameter optimisation to obtain unbiased estimates
on this small data. Consistent with expectations for limited
tabular data, the gradient-boosted tree methods outperformed the linear,
kernel, and neural alternatives. 

Theere are around 50 features in total, and would be adding it soon to the paper along with their importance.
The ablation study must also be done to see the contribution of each feature to the model performance.

[TODO: list the exact algorithms
and their key settings in a table and explain the hyperparameter tuning strategy in more detail along with the features and their importance.]

%==============================================================
\section{Metrics}\label{sec:metrics}
%==============================================================
Performance was assessed with five-fold cross-validation, reporting
accuracy and the F1 score for GOF detection given the class imbalance.
The strongest models reached an F1 of roughly 0.67 at about 75\% accuracy,
well above a random baseline near 0.38, indicating genuine signal in the
engineered features. Adding the structural descriptors produced a modest
but consistent improvement over sequence-only features. 

[TODO: add a
compact results table and confidence intervals.]

%==============================================================
\section{Conclusion}\label{sec:conclusion}
%==============================================================
These results indicate that direction-aware, receptor-specific prediction
of nAChR variant effects is feasible: engineered sequence and structural
features carry enough signal to separate LOF from GOF above chance. Such a
predictor could triage the expanding population of nAChR VUS, prioritising
which variants warrant costly functional follow-up and supporting
precision therapy. The main limitations are data scarcity and the
treatment of LOF and GOF as a clean dichotomy; future work should expand
the curated dataset and establish a shared direction-of-effect benchmark
for nAChRs.

%==============================================================
% References
%==============================================================
\begin{thebibliography}{9}

\bibitem{sift} P.~C. Ng and S.~Henikoff, SIFT: predicting amino acid
changes that affect protein function, {\em Nucleic Acids Res.}
{\bf 31}, 3812 (2003).

\bibitem{polyphen2} I.~A. Adzhubei {\em et al.}, A method and server
for predicting damaging missense mutations, {\em Nat. Methods}
{\bf 7}, 248 (2010).

\bibitem{cadd} M.~Kircher {\em et al.}, A general framework for
estimating the relative pathogenicity of human genetic variants,
{\em Nat. Genet.} {\bf 46}, 310 (2014).

\bibitem{alphamissense} J.~Cheng {\em et al.}, Accurate proteome-wide
missense variant effect prediction with AlphaMissense, {\em Science}
{\bf 381}, eadg7492 (2023).

\bibitem{esm1b} N.~Brandes {\em et al.}, Genome-wide prediction of
disease variant effects with a deep protein language model,
{\em Nat. Genet.} {\bf 55}, 1512 (2023).

\bibitem{funncion} A.~Brunklaus {\em et al.}, Gene variant effects
across sodium channelopathies predict function and guide precision
therapy, {\em Brain} {\bf 145}, 4275 (2022).

\bibitem{logofunc} D.~Stein {\em et al.}, Genome-wide prediction of
pathogenic gain- and loss-of-function variants from ensemble learning
of a diverse feature set, {\em Genome Med.} {\bf 15}, 103 (2023).

\bibitem{premode} G.~Zhang {\em et al.}, Predicting the mode of action
of missense variants (PreMode), bioRxiv (2024).

\bibitem{mohanraj2026} H.~K. Mohan Raj {\em et al.}, A curated dataset of
experimentally characterised missense variants in human nicotinic
acetylcholine receptors, (2026).


\end{thebibliography}

\end{document}
